% ---------------------------------------------------------------------------
% JudgeSense --- TMLR submission.
%
% TMLR ships its own style file (tmlr.sty + tmlr.bst) from the author kit at
% https://jmlr.org/tmlr/. Drop those beside this file and swap the two marked
% lines below; the body needs no other change. Until then this builds with
% article so the document is checkable without the kit.
% ---------------------------------------------------------------------------
\documentclass[11pt]{article}          % TMLR: \documentclass{article} + \usepackage{tmlr}
\usepackage[margin=1in]{geometry}      % TMLR: delete (tmlr.sty sets the page)

\usepackage[utf8]{inputenc}
\usepackage[T1]{fontenc}
\usepackage{lmodern}                   % scalable fonts; microtype requires them
\usepackage{booktabs}
\usepackage{amsmath,amssymb}
\usepackage{graphicx}
\usepackage{microtype}
\usepackage[numbers,round]{natbib}
\usepackage[hidelinks]{hyperref}
\usepackage{url}

\title{JudgeSense: Measuring the Prompt Sensitivity of LLM-as-a-Judge}

% TMLR review copies are anonymous; the kit's \author block handles this.
\author{Rohith Reddy Bellibatlu\thanks{Affiliation and contact to be completed
before submission.}}
\date{}

\begin{document}
\maketitle

\begin{abstract}
Language models are routinely used to evaluate other language models, and an
evaluator earns that role by being reproducible. We ask a question the usual
validation cannot answer: does a judge return the same verdict when the same
request is worded differently? We introduce JudgeSense, a benchmark of 880 items
across four evaluation tasks drawn from human-labelled corpora, each issued under
two instruction templates that differ in wording and not in what they ask. The
metric never consults the ground-truth label --- a judge that is uniformly wrong
but uniformly consistent scores perfectly --- and it is reported against each
judge's own repeat-agreement ceiling, so ordinary decoding variance is separated
from sensitivity to wording rather than attributed to it. We release the
heuristic baselines that would beat each task, with intervals and support counts,
alongside a pre-registered floor that a competent judge must clear for a task to
count as discriminating at all. Across three judges, meaning-preserving rewording
costs real agreement in every measurable cell, and the size of that loss depends
more on the evaluation task than on the model performing it: the ordering of
tasks by sensitivity is identical across judges, and the spread within a task is
smaller than the spread within a judge across tasks. We also report a
judge-specific refusal behaviour that a naive harness would misreport as a
format-adherence failure, and the construction defects we found in our own
benchmark, with the measurements that exposed them.
\end{abstract}

\section{Introduction}
\label{sec:intro}

Language models are now routinely used to evaluate the output of other language
models. The practice is cheap, fast, and correlates well enough with human
preference to have displaced human annotation across large parts of the
evaluation stack \citep{zheng2023judging}. Leaderboards are built on it, training
signals are derived from it, and papers report it as evidence.

An evaluator earns that role by being reproducible. If two researchers ask the
same judge the same question in different words and receive different verdicts,
the judge is not measuring the property it claims to measure. It is partly
measuring the phrasing. That failure is invisible in the usual validation, which
asks whether a judge agrees with humans on average. A judge can correlate well
with human ratings in aggregate while being unstable on any individual item, and
aggregate correlation cannot distinguish the two.

Instability under rephrasing is well documented for language models performing
ordinary tasks. Small, meaning-preserving changes to a prompt (formatting, ordering, the choice
between synonyms) move accuracy substantially
\citep{sclar2024quantifying, mizrahi2024state, elazar2021measuring}. The
evaluation literature has separately established that judges carry systematic
biases: they favour the first-presented candidate
\citep{wang2023fair, shi2024judging}, and length is a well-known confound --- in
the human preference votes we draw on, annotators chose the longer response in
roughly 70\% of usable comparisons. What has been measured less is the
property that matters most for a measuring instrument: whether the same judge,
asked the same question two ways, returns the same answer.

This paper measures that directly. We introduce JudgeSense, a benchmark of 880
items across four evaluation tasks, each issued under two instruction templates
that differ in wording and not in what they ask. The central quantity is not
accuracy. A judge that is wrong on every item but wrong \emph{identically} under
both phrasings is perfectly reproducible, and our metric assigns it a perfect
score. We are asking whether the instrument is stable, and treating whether it is
correct as a separate question.

\paragraph{Separating sensitivity from noise.}
A judge sampling at nonzero temperature disagrees with itself on the identical
prompt, so some part of any measured disagreement is ordinary decoding variance
rather than sensitivity to wording. Reporting raw agreement attributes all of it
to paraphrasing. We therefore issue each template a second time, unchanged, and
report the \emph{difference} between paraphrase agreement and repeat agreement.
The repeat arm is the judge's own reproducibility ceiling on that task, and the
gap below it is the quantity the paper is about. This distinction is not
cosmetic: in a pilot we discarded, a misconfigured decoding parameter produced a
repeat ceiling low enough to absorb the entire effect we report, and nothing in
the output distinguished it from a genuine one.

\paragraph{A benchmark has to be attackable before it is trustworthy.}
A measurement of judge stability is only as good as the items it is measured on.
If a task can be passed by a heuristic that ignores the intended construct, then
a high score says nothing about judgment. Always answer A. Pick the longer
candidate. Pick the passage sharing more words with the query. We therefore
run those heuristics against the released files and report what they score, with
intervals and support counts. None beats chance, and we report which control is
weakest and why. Equally, a benchmark on which \emph{nothing} works,
including a competent judge, is shortcut-free and useless; so we also declare a
threshold the best judge must clear for a task to count as discriminating at
all, and report each task against the rate a constant-answer judge achieves
rather than against uniform chance. Both directions are checked, because a construction that
suppresses exploitable signal can suppress genuine signal by the same mechanism.

\paragraph{This paper is a correction as well as a contribution.}
% The self-citation names all three authors in the bibliography, and the
% surrounding sentence claims the work as ours, so rendering it under review
% would deanonymise the submission as surely as the author block would.
An earlier version of this benchmark
\ifblind(reference withheld for review)\else\citep{bellibatlu2026judgesense}\fi{} was
withdrawn after review. Its items were
hardcoded in source while being described as benchmark-sourced; it shipped 75
unique items; it presented the correct candidate first in every pairwise item, so
a judge that always answered ``A'' was indistinguishable from a perfect one; and
one factuality template inverted the answer polarity, inflating measured
disagreement uniformly across every judge. We rebuilt from real corpora with
per-item provenance and fixed each defect. In doing so we found further defects of the same class in our own rebuild: a
template assignment that functioned as an answer key, candidate responses
carrying contradictory gold labels, and a length-balancing step that installed a
difficulty confound in place of the verbosity one it removed. We report those
too, with the measurements that exposed them. The failure mode we are documenting is not exotic. It is
ordinary, it is quiet, and it survives every check that does not specifically
look for it.

\paragraph{Contributions.}
\begin{enumerate}
  \item \textbf{A benchmark with attack surface measured, not asserted.} 880
    items from four human-labelled corpora with per-item provenance, released
    with the heuristic baselines that would beat it, their confidence intervals,
    and their support counts.
  \item \textbf{An endpoint measured against each judge's own noise ceiling.}
    $\Delta\mathrm{JSS}$, the gap between paraphrase agreement and a judge's own
    repeat-agreement ceiling. This demotes decoding variance rather than
    eliminating it, and we state the estimator's identity so a reader can see
    what it does and does not isolate. The endpoint, the estimator and the
    item-level clustering unit were pre-registered before the sweep; the support
    floors and the smallest effect of interest were added after the first
    judge's results and are reported as declared rather than pre-registered.
  \item \textbf{Refusal treated as a measurement rather than a nuisance.}
    Provider-flagged declines are recorded and reported separately from
    unparseable output, which are indistinguishable after parsing but are
    different facts about a judge.
  \item \textbf{Evidence that the ordinal task is the least stable.} Coherence
    carries the largest effect for all three judges and factuality the smallest,
    and on coherence the variation across judges is smaller than the variation
    across tasks within a judge. All three judges come from one vendor family, so
    this is a within-family observation and we do not claim it generalises.
\end{enumerate}

\section{Related work}
\label{sec:related}

\paragraph{LLM judges and their validation.}
Using a strong language model to score the output of another is now standard
practice, established at scale by MT-Bench and Chatbot Arena
\citep{zheng2023judging} and adopted across holistic evaluation efforts
\citep{liang2023holistic}. Validation of these judges is almost always
\emph{correlational}: a judge is trusted when its scores agree with human ratings
in aggregate \citep{chiang2023can, liu2023geval}. That criterion is necessary and
not sufficient. Aggregate correlation is compatible with per-item instability,
because errors that cancel across a corpus still make any individual verdict
unreliable, and a researcher who reruns an evaluation with differently worded
instructions has no guarantee of reproducing it. Our measurement is orthogonal to
accuracy by construction: it never consults the ground-truth label, and a judge
that is uniformly wrong scores perfectly on it.

\paragraph{Known judge biases.}
A body of work documents specific, systematic ways judges fail. Position bias ---
preferring whichever candidate is presented first --- is strong enough to require
explicit correction \citep{wang2023fair} and has been characterised
systematically across pairwise comparative assessment \citep{shi2024judging}.
Length is the other well-known confound, and rather than lean on a citation we
measure it directly in our own source data: in the MT-Bench human preference
votes we draw from, annotators preferred the longer of the two responses in
roughly 70\% of usable comparisons (Section~\ref{sec:shortcuts}), so a judge that
simply picked the longer answer would score far above chance on an unfiltered
sample of them.

These are findings about \emph{what} judges prefer. We take them as constraints
on benchmark construction rather than as our object of study: because position
and length bias are known, a benchmark that does not control them measures the
bias instead of the judge. We therefore present both candidate orderings for
every pairwise item and balance the length distribution exactly, then verify on
the released files that neither heuristic beats chance.

\paragraph{Prompt sensitivity in language models generally.}
That model behaviour moves under meaning-preserving reformulation is
well-established outside the evaluation setting. Formatting choices alone shift
accuracy across a wide spread \citep{sclar2024quantifying}; paraphrasing task
instructions changes measured capability enough to alter model rankings
\citep{mizrahi2024state}; consistency under rephrasing has been treated as a
first-class property of pretrained models \citep{elazar2021measuring}; and the
sensitivity itself has been assessed and benchmarked as a phenomenon in its own
right \citep{zhuo2024prosa, razavi2025benchmarking}. Systematic surveys of
evaluation practice reach the same conclusion from the other direction
\citep{wei2024systematic}. All of this concerns a model \emph{being tested}. Our
contribution is to apply the same lens to the model doing the testing, where
instability propagates silently into every downstream number.

\paragraph{Judging judges.}
Recent work evaluates LLM judges as objects in their own right, including
alignment with human labels \citep{thakur2024judging} and general capability
benchmarks \citep{li2024llms}. These efforts mostly ask whether judges are
\emph{right}.

The closest prior work to ours is \citet{arabzadeh2025human}, who study prompt
sensitivity in LLM-based relevance judgment directly, comparing human- and
LLM-authored prompt variants across binary, graded and pairwise formulations of
the same relevance task. The difference is in what the measurement is referenced
against. That work evaluates labels produced under varying prompts against human
relevance gold, so it measures how rewording moves a judge's \emph{accuracy};
JSS never consults the gold label and measures how rewording moves a judge's
\emph{agreement with itself}. The two are complementary and can diverge: a judge
whose accuracy is unchanged under rewording may still be flipping individual
verdicts in compensating directions, which our measurement sees and an
accuracy-referenced one does not. We also separate that quantity from decoding
variance by measuring each judge against its own repeat-agreement ceiling rather
than against an assumed noise floor of zero.

\paragraph{Sources and statistical machinery.}
The items come from existing human-labelled corpora: TruthfulQA for factual
accuracy \citep{lin2022truthfulqa}, SummEval for expert coherence ratings
\citep{fabbri2021summeval}, BEIR's TREC-COVID split for graded relevance
judgements \citep{thakur2021beir}, and MT-Bench's human preference votes
\citep{zheng2023judging}. No label in this benchmark is assigned or inferred by
us. Agreement is chance-corrected with Cohen's $\kappa$ \citep{cohen1960coefficient},
quadratic-weighted on the ordinal task, and all intervals are cluster bootstrap
\citep{efron1979bootstrap} resampled at the item, since the two candidate
orderings of a pairwise item and its repeat arms are not independent
observations. The judges evaluated are drawn from current instruction-tuned and
reasoning-tuned families \citep{openai2023gpt4, gemini2023gemini,
grattafiori2024llama3, qwen2024qwen25, jiang2023mistral, deepseek2025r1}.

% ---------------------------------------------------------------------------
% Benchmark construction. Number-independent: everything here is a property of
% the artifact, not of any judge's behaviour, so it is stable across runs.
% ---------------------------------------------------------------------------
\section{The JudgeSense benchmark}
\label{sec:benchmark}

JudgeSense measures whether an LLM judge returns the same verdict when the same
evaluation request is worded differently. The construct is deliberately narrow:
we are not asking whether a judge is \emph{correct}, but whether it is
\emph{stable}. A judge that is consistently wrong is reproducible; a judge whose
verdict depends on which of two equivalent phrasings it received is not usable as
a measuring instrument, however accurate it looks on average.

\subsection{Tasks and sources}
\label{sec:sources}

The benchmark covers four evaluation tasks, each drawn from an existing
human-labelled corpus. No item is authored by us, and no label is assigned or
inferred by us; every label is the one the source corpus already carried.

\begin{description}
  \item[Factuality] (250 items, TruthfulQA). A statement is judged accurate or
    inaccurate. Items whose question is ill-posed, and whose only defensible
    answer is a refusal to answer, are excluded at load time.
  \item[Coherence] (250 items, SummEval). A summary is rated 1--5 for
    coherence against expert annotations. Summaries with identical text but
    conflicting expert labels are dropped rather than merged.
  \item[Relevance] (250 items, 500 rows, BEIR TREC-COVID). Given a topic, the
    judge picks the more relevant of two passages. \emph{Both} candidates carry
    a human relevance grade: the positive is graded fully relevant (2) and the
    distractor is one a human assessor explicitly rejected (0). Partially
    relevant documents (grade 1) are used for neither role.
  \item[Preference] (130 items, 260 rows, MT-Bench human judgments). The judge
    picks the response humans preferred. Only comparisons with at least two
    annotator votes and a clear majority are used.
\end{description}

The relevance construction is the substantive departure from the usual practice
of drawing distractors from a positives-only collection. There, a negative can
only be a document \emph{absent} from an incomplete relevance judgment set, so a
topically similar distractor may in fact be relevant but unjudged, and the item
penalises a defensible answer. TREC-COVID carries graded judgments including
tens of thousands of explicit non-relevant assessments, so here both candidates
are human-labelled and the distractor is one a human affirmatively rejected.

Every shipped item records its source dataset, split, per-record identifier, and
the source configuration where the source defines one, so any item can be traced
to the row it came from.
Loaders raise rather than substitute when a source is unavailable; there is no
fallback path that could silently generate an item.

\subsection{Paraphrase design}

Each item is issued under two instruction templates that differ in wording and
not in what they ask. Templates are held to a fixed label space: a check run in
continuous integration asserts that no factuality or coherence template alters
the set of admissible answers or inverts their polarity. The pairwise templates
share a two-option answer space by construction and are not covered by that
check.

This constraint is a correction of our own earlier design. An earlier version of
this benchmark included a factuality template that inverted the answer polarity
(``does this response contain factual errors? answer NO if accurate''). Under a
single shared label space, ``YES'' then denoted opposite conclusions in different
arms, so a judge that answered both arms correctly was recorded as having
disagreed with itself. That is a label-convention defect in the benchmark rather
than a property of any model, and it measures a different construct: handling a
polarity flip is a comprehension task, whereas we ask whether equivalent
phrasings yield equivalent verdicts.

We state the direction of the resulting bias rather than leave it implicit.
Polarity-inverted pairings are among the hardest items in the set, so excluding
them can only push measured consistency \emph{upward}; our stability figures are
conditioned on that exclusion and should be read accordingly. The alternative ---
retaining the inverted templates and remapping each template's raw answers onto a
single canonical decision space before comparison --- is implemented in the
released code, and we report the remapped analysis alongside the
excluded-pairing analysis so the exclusion's effect on the reported numbers is
visible rather than assumed.

\subsection{Position control}

Both pairwise tasks present every item twice, once in each candidate ordering.
The correct answer therefore sits at position A in exactly half the rows, and a
judge that ignores content and always answers ``A'' scores $0.500$ rather than
$1.000$.

This too corrects an earlier version, in which the correct candidate always
appeared first. Under that construction a position-anchored judge was
indistinguishable from a perfect one, and most judges scored a degenerate
$\mathrm{JSS}=1.000$ that reflected the layout rather than any judgment.

\subsection{Shortcut controls}
\label{sec:shortcuts}

A benchmark measures what it claims to measure only if it cannot be passed by a
heuristic that ignores the intended construct. We therefore evaluate the shipped
splits against judges that use no judgment at all. Each is run over the released
files, not over the loaders, so these are properties of the artifact:

\begin{itemize}
  \item \textbf{Position only} (always answer A): $0.500$ on relevance
    ($n=500$) and $0.500$ on preference ($n=260$), by construction of the swap
    design.
  \item \textbf{Length only} (pick the longer candidate): $0.482$ on relevance
    ($n=494$), $0.508$ on preference ($n=256$).
  \item \textbf{Lexical overlap only} (pick the candidate sharing more terms with
    the query): $0.540$ on relevance ($n=252$).
\end{itemize}

Support counts fall below the row counts because a heuristic that ties has no
signal to act on; tied rows are excluded rather than split, which is why the
overlap control rests on roughly half the relevance rows. All three controls are
computed from the released files by \texttt{scripts/shortcut\_controls.py}, whose
output is committed, so every figure above is reproducible without rerunning any
judge.

None of these exceeds chance in either direction, though the overlap control is
the weakest of the three and we do not overstate it: its 95\% Wilson interval is
$[0.478, 0.600]$, so the data are consistent with a residual lexical signal worth
about $0.57$. It also reads only half the relevance rows, because the remainder
tie on the overlap statistic and a heuristic with no signal on a row is not
evidence about that row. The length and position controls are far tighter.

The length control is enforced rather than observed. Human annotators in the
preference corpus preferred the longer response in roughly 70\% of usable
comparisons, so an unfiltered sample rewards verbosity directly; we hold the
winner-longer share at exactly 50\% by taking an equal number of items from each
length bucket. That $50\%$ is a construction invariant of the item set and is a
different quantity from the $0.508$ above, which is what a length-following judge
scores on the rendered rows once ties are dropped.

The preference split ships 130 items rather than the round 250, and we state the
full attribution rather than the net gap. Of the source comparisons, 607 failed
the decisive-vote label rule or carried no strict majority, one was duplicate
content, 67 were dropped because their gold contradicted another item's gold on
shared candidate text, and 102 more were dropped to hold the length balance
exactly. Each filter removes a way the benchmark could be passed or contradicted
without judging, and the last two were added after we found both failures in our
own shipped data. We regard the reduced split as the correct trade: the
alternative is a larger set that a verbosity heuristic beats, and one that scores
a consistent judge as wrong whenever the same response appears on both sides of
two different comparisons.

Balancing on length alone would install a second confound in place of the first.
Because the winner-shorter bucket ships whole and only the winner-longer bucket
is cut, any criterion used to make that cut applies to one bucket and not the
other. Cutting by vote margin --- keeping the most contested comparisons, as an
earlier build did --- left the buckets matched on length but far apart on
annotator agreement: mean vote-margin ratio $0.567$ against $0.745$, and $16\%$
against $54\%$ unanimous. A judge that is stronger on close calls would then
read as length-biased. We therefore \emph{stratify}: the winner-shorter bucket
fixes the target distribution of vote shapes (absolute margin and total votes),
and the winner-longer bucket is filled stratum by stratum to match it. The
shipped buckets agree to $0.0003$ in mean vote-margin ratio ($0.8603$ against
$0.8605$), have identical medians, and contain the same number of unanimous
comparisons ($49$ each).

Matching stops the split from collapsing onto coin-flips, but the opposite
failure is equally real and less obvious. The filter that removes contradictory
gold is a greedy selection over a conflict graph, so the order in which items are
considered determines both which survive and how hard the survivors are. Keeping
the firmest human majorities first is the intuitive rule and is worse on both
counts, measured on this pool: $106$ items at $90.6\%$ unanimous, against $130$
items at $75.4\%$ when the most contested are kept first. A split on which the
annotators were nearly always unanimous cannot separate a competent judge from a
mediocre one --- both score near the top --- which is the ceiling effect that
made the earlier version of this benchmark uninformative. We therefore resolve
conflicts contested-first, and report the resulting distribution rather than only
the balance: $75.4\%$ of shipped comparisons carry a unanimous human majority and
the rest are genuine close calls.

The lexical control required a similar correction. Selecting each distractor to
maximise its term overlap with the topic --- the obvious way to make an item
``hard'' --- made distractors \emph{denser} in query terms than the genuinely
relevant passages, so the heuristic ``pick the passage with \emph{lower} overlap''
recovered the answer far above chance. Matching the distractor's retrieval score
to the positive's, rather than maximising it, removes overlap as a signal in
either direction. We report the control two-sided for this reason: a split whose
answer is reliably the \emph{shorter} or \emph{less overlapping} candidate is
exploited exactly as cheaply as one biased the other way.

\subsection{Release and integrity}

The dataset is released with Croissant metadata, per-item provenance, and a
content hash for each split that excludes retrieval timestamps, so a rebuild can
be verified as identical in content without being byte-identical in metadata. A
28-check audit gate runs in continuous integration and fails the build on
duplicate content, unbacked provenance claims, label degeneracy, insufficient
clustering support, contradictory ground truth, or a length shortcut exceeding
its bound. Checks that cannot read the data they are meant to inspect fail
rather than pass, so a check that silently measures nothing cannot be mistaken
for a check that found nothing wrong.

% ---------------------------------------------------------------------------
% Metric definitions and the analysis plan. Number-independent.
% ---------------------------------------------------------------------------
\section{Measuring sensitivity}
\label{sec:metrics}

\subsection{Judge Sensitivity Score}

For an item issued under two equivalent templates, let $d_1(p)$ and $d_2(p)$ be
the judge's parsed decisions on pair $p$. The Judge Sensitivity Score is the
agreement rate

\[
  \mathrm{JSS} = \frac{1}{|P|}\sum_{p \in P} \mathbb{1}[d_1(p) = d_2(p)],
\]

over the set of scored pairs $P$. JSS never consults the ground-truth label. A
judge that is wrong on every item but wrong \emph{identically} under both
phrasings scores $1.0$, and this is intended: it is a measure of reproducibility,
not of accuracy. We report accuracy separately, and the two answer different
questions. On the pairwise tasks the companion measure is \emph{position-corrected}
accuracy --- whether the judge selected the side carrying the ground truth, given
that the swap design places it at position A in exactly half the rows --- and on
the pointwise tasks it is ordinary accuracy against the source label. The
distinction matters: a judge that always answers ``A'' has an ordinary accuracy
that depends only on the layout, while its position-corrected accuracy is
$0.500$ and its behaviour is immediately visible.

Because raw agreement is inflated when a judge's output distribution is
compressed --- a judge answering ``YES'' almost always agrees with itself --- we
report the chance-corrected form (Cohen's $\kappa$; quadratic-weighted for the
ordinal coherence task) alongside every raw figure.

\subsection{The repeat baseline, and why it is the endpoint}
\label{sec:repeat}

Raw JSS below $1.0$ does not by itself demonstrate sensitivity to \emph{wording}.
A judge sampling at nonzero temperature, or one served through a
non-deterministic backend, disagrees with itself on the identical prompt. Some
portion of any measured disagreement is therefore ordinary decoding noise, and a
paper that reports raw JSS attributes all of it to paraphrasing.

We separate the two by issuing one template a second time, unchanged. The
resulting $\mathrm{JSS}_{\text{rep}}$ is that judge's reproducibility ceiling on
that task: the agreement it achieves when nothing varies. The primary endpoint is
the difference

\[
  \Delta\mathrm{JSS} = \mathrm{JSS}_{\text{paraphrase}} - \mathrm{JSS}_{\text{rep}},
\]

which is negative when rewording costs more stability than resampling does. This
is the quantity the paper's claim is about, and it was fixed in a pre-registered
analysis plan committed to the repository before any judge was run.

\subsection{Uncertainty}

All intervals are 95\% cluster bootstrap intervals over 2{,}000 resamples, with
the item as the resampling unit. Item-level clustering is required rather than
preferred: a pairwise item contributes two rows (one per candidate ordering) and
its repeat arm nests inside it, so resampling rows independently treats
correlated observations as independent and understates uncertainty. The
clustering unit is a mandatory argument in our implementation, with no default,
so it cannot be omitted by accident.

An effect is reported as present for a judge--task cell when the interval for
$\Delta\mathrm{JSS}$ excludes zero. The direction is read from the sign.

\subsection{Outcome categories: verdict, refusal, malformed}
\label{sec:refusal}

Each call terminates in exactly one of three outcomes.

A \textbf{verdict} is a response the strict parser maps to a label in the task's
answer space. A \textbf{refusal} is a response the provider itself flags as
declined, typically with empty content. Providers spell this differently ---
\texttt{stop\_reason="refusal"} for Anthropic, \texttt{finish\_reason=
"content\_filter"} for OpenAI-compatible APIs, \texttt{SAFETY} or
\texttt{RECITATION} for Google --- so the accepted set is enumerated in a single
place in the released code and matched case-insensitively, and the raw string the
provider returned is retained on every record. Recognising only one family's
spelling would report a refusal rate of zero for every other family regardless of
what those judges did, and would silently charge their declines as paraphrase
disagreement.
\textbf{Malformed output} is a completed response the parser cannot map. Parsing
is strict and never coerces: an unmappable response is recorded with its raw text
retained, never rounded to the nearest label.

Malformed output and refusal are treated differently, and the asymmetry is
deliberate. A malformed arm is \emph{charged as disagreement} and stays in the
support: the model completed a response, so the pair reflects a judgment the
parser could not read rather than one that was never rendered. A refused arm is
removed from the support, because no judgment was rendered at all. Malformed-output
rate is reported per cell alongside the refusal statistics below, so neither is
absorbed into the other.

Refusals and malformed output are indistinguishable after parsing --- both yield
no usable label --- but they are different measurements, and we do not merge
them. A refusal occurs \emph{upstream} of any judgment: the provider halted
before the model rendered a verdict. Scoring it as paraphrase disagreement
asserts that the judge produced two conflicting judgments, which it did not.
Treating refusal as a third label is worse still, since a judge that refuses
everything would then score $\mathrm{JSS} = 1.0$.

We therefore compute JSS over pairs in which both arms returned verdicts, and
report the support size $n$ alongside every figure. Refusal is reported as its
own construct, per cell:

\begin{itemize}
  \item \emph{refusal rate}: refused calls over all calls;
  \item \emph{refusal discordance rate} (RDR): pairs in which exactly one arm was
    refused;
  \item \emph{consistent refusal rate}: pairs in which both arms were refused.
\end{itemize}

RDR is itself a sensitivity statistic, and the most interesting of the three: a
nonzero RDR means a meaning-preserving rewording changed whether the judge was
willing to judge at all. That is prompt sensitivity of the willingness-to-answer
layer, and it is invisible to any analysis that scores refusals as disagreement
or discards them.

Conditioning on verdict pairs restricts a refusing judge to the items it agreed
to judge, which raises the fair objection that judges are then compared on
different subsets. We answer it with a measurement rather than an argument: every
cross-judge comparison in a task with nonzero refusals is recomputed on the
common support --- the items every compared judge answered under both phrasings
--- and we additionally report a refusal-inclusive figure in which every refused
arm counts as disagreement, the most punitive available reading. Both appear
beside the conditioned figure, so the effect of the conditioning is visible
rather than asserted.

\subsection{Decoding budget}
\label{sec:budget}

Every judge is run with an identical output token limit. This is not a neutral
default: a class-dependent limit --- a tight cap for instruction-tuned judges and
a generous one for reasoning-tuned judges, on the reasoning that the latter emit
chain-of-thought --- confounds judge class with inference configuration, so any
difference between the classes admits an explanation that has nothing to do with
the models. It also truncates real responses: under such a configuration we
observed judges terminating at the cap mid-sentence, with the fragment retained
as the answer. Disagreement so produced is a property of the cap rather than of
the judge. We do not quote a rate here, because the rate is a property of the
configuration under study and is reported with its denominator by the ablation
below rather than asserted in passing.

The matched budget removes the explanation by construction. Because a token
limit is a ceiling and not a target, compliant judges that emit a single-token
answer are unaffected in cost; only verbose non-compliers cost more, and those
are precisely the calls that were previously being truncated into noise. We
quantify the budget's effect directly rather than assert it away, reporting a
subsample ablation under the class-dependent configuration together with the
cap-termination rate under each (Appendix~\ref{app:budget}).

\section{Results}
\label{sec:results}

We report three judges from one vendor family --- Claude Haiku 4.5, Sonnet 4.5
and Opus 4.7 --- on the four tasks. Factuality, coherence and relevance are
complete for all three; the preference split is complete for Haiku only, and
Section~\ref{sec:results-limits} states what that costs. Every figure below is
produced by \texttt{scripts/regenerate\_results.py} from the committed raw
outputs; none is transcribed.

\subsection{Paraphrasing costs stability beyond decoding noise}

Table~\ref{tab:main} reports $\Delta\mathrm{JSS} = \mathrm{JSS}_{\text{para}} -
\mathrm{JSS}_{\text{rep}}$, the pre-registered endpoint. In all ten cells with
sufficient support the interval excludes zero and the effect is negative: a
meaning-preserving rewording costs more agreement than re-issuing the identical
prompt does.

The magnitudes are not uniform. They span an order of magnitude, from $-0.030$
(Opus, factuality) to $-0.206$ (Haiku, coherence), and all exceed the
$|\Delta\mathrm{JSS}| < 0.02$ threshold declared in advance as not practically
meaningful.

The repeat ceiling itself is high everywhere: $\mathrm{JSS}_{\text{rep}}$ lies
between $0.916$ and $1.000$ across the ten cells. At a matched decoding budget
and temperature zero these judges are close to deterministic on a fixed prompt,
which is what makes the paraphrase gap interpretable. It is also what a
mismatched configuration destroys: in a discarded pilot in which one vendor's
calls omitted the temperature parameter and defaulted to $1.0$, the same judge
and task returned $\mathrm{JSS}_{\text{rep}} = 0.864$ --- a $13.6\%$
self-disagreement rate on byte-identical prompts, large enough to absorb the
entire effect reported here. We note this because the failure is silent: nothing
in the output distinguishes a decoding-noise ceiling from a misconfigured one.

\subsection{Sensitivity tracks the task more than the judge}

The ordering of tasks by sensitivity is the same for every judge. Coherence is
the most sensitive task in all three cases, factuality the least, with the
pairwise tasks between:

\begin{center}
\begin{tabular}{lcccc}
\toprule
Judge & coherence & preference & relevance & factuality \
\midrule
Haiku 4.5   & $-0.206$ & $-0.081$ & $-0.058$ & $-0.044$ \
Sonnet 4.5  & $-0.162$ & ---      & $-0.158$ & $-0.038$ \
Opus 4.7    & $-0.180$ & ---      & $-0.036$ & $-0.030$ \
\bottomrule
\end{tabular}
\end{center}

Within a task the spread across judges is smaller than the spread across tasks
within a judge. Coherence ranges $-0.162$ to $-0.206$ across three judges, while
Haiku alone ranges $-0.044$ to $-0.206$ across four tasks. On this evidence,
which evaluation you are asking a judge to perform predicts paraphrase
sensitivity better than which judge you asked.

The exception is Sonnet on relevance ($-0.158$ against $-0.058$ and $-0.036$ for
the other two), and it has a specific cause treated in
Section~\ref{sec:results-refusal} rather than reflecting sensitivity to wording.

Coherence being the most sensitive task is consistent with its structure. It is
the only ordinal task, so a judge has five admissible answers rather than two,
and a rewording that shifts a judgement by one scale point counts as
disagreement. Chance-corrected agreement supports this reading: on the quadratic
weighted kappa, which charges an adjacent-category disagreement less than a
distant one, coherence scores $0.859$--$0.874$, far above its unweighted kappa of
$0.594$--$0.727$. Most coherence disagreements are one-point moves.

\subsection{The tasks discriminate}
\label{sec:results-discriminate}

A benchmark on which no heuristic beats chance is worthless if no competent judge
beats chance either. Section~\ref{sec:shortcuts} established the floor; the
pre-registered ceiling requires the best judge on each task to exceed a
threshold declared before the run. All four clear it:

\begin{center}
\begin{tabular}{lccc}
\toprule
Task & required & best observed & chance \
\midrule
factuality & $0.75$ & $0.944$ (Opus)   & $0.50$ \
coherence  & $0.40$ & $0.452$ (Sonnet) & $0.20$ \
relevance  & $0.70$ & $0.882$ (Opus)   & $0.50$ \
preference & $0.65$ & $0.873$ (Haiku)  & $0.50$ \
\bottomrule
\end{tabular}
\end{center}

Coherence clears its threshold by the narrowest margin, which is expected on a
five-point scale whose gold labels are a rounded three-annotator mean; the
weighted kappa above is the better measure of that task's signal.

Position-corrected accuracy also confirms the swap design is working on live
judges rather than only on the artifact. Answer-A rates are $0.494$--$0.569$ on
relevance and $0.500$--$0.516$ on preference, against the degenerate $1.000$ that
the previous version of this benchmark produced for eleven of thirteen judges.

\subsection{Refusal is a distinct failure mode, and it is judge-specific}
\label{sec:results-refusal}

Sonnet declined $24.9\%$ of relevance arm calls, reported by the provider as a
refusal with empty content. Opus refused $5.2\%$ of the same prompts and Haiku
none at all. The corpus is TREC-COVID, so the items are medical, and the effect
is a property of one judge's safety behaviour rather than of the benchmark: the
three judges saw identical prompts.

Had refusals been folded into malformed output, Sonnet's relevance cell would
report a $22.4\%$ format-adherence failure. It is not one. Separating the two
outcome categories is what makes the reported malformed rate a statement about
instruction-following, and the refusal rate a statement about willingness to
judge.

The refusal-conditioned JSS is identified only within bounds, because a rewording
can itself change whether the judge answers --- refusal is an outcome, not a
nuisance. We therefore report the worst-case interval alongside the point
estimate wherever refusals occur, rather than presenting the conditioned value
alone.

\subsection{What this evidence does not cover}
\label{sec:results-limits}

Three judges, one vendor. Every result here is within the Claude family, so the
consistency of the task ordering across judges is a within-family observation and
may not extend to other vendors. The registry and harness support fourteen
judges across seven providers, and the remaining eleven are unrun.

The preference split is complete for Haiku only. Sonnet and Opus stopped short at
$44$ and $170$ of $260$ rows when the API balance was exhausted mid-run; the
pipeline refuses to emit $\Delta\mathrm{JSS}$ below the pre-registered floor of
100 clusters, so those two cells report no endpoint rather than one computed
from $17$ and $80$ clusters. The rows already collected are retained.

One judge is not decoding-matched. Opus 4.7 rejects an explicit temperature
parameter, so it ran at the provider default while the other two ran at zero.
This is recorded on every one of its call records rather than inferred, and its
cells should be read with that in mind: its ceiling of $0.916$ on coherence is
the lowest observed, which is the direction an unmatched temperature would push.

Finally, the design measures sensitivity within an item across two templates. It
cannot compare templates to one another, because no item appears under more than
one template pair (Section~\ref{sec:benchmark}).

\section{Discussion}
\label{sec:discussion}

\subsection{What the result means for practice}

The finding that sensitivity tracks the task more than the judge has a direct
consequence: choosing a stronger judge is not a reliable way to buy stability.
Across the three judges measured, the most and least sensitive judge on coherence
differ by $0.044$ $\Delta\mathrm{JSS}$, while the most and least sensitive
\emph{task} within a single judge differ by $0.162$. A practitioner worried about
reproducibility gains more by reconsidering what they are asking the judge to do
than by upgrading the model.

The ordinal task is the least stable of the four, and we think the mechanism is
general rather than particular to coherence. A five-point scale gives a rewording
more ways to move the answer, and the weighted kappa shows most of that movement
is a single scale point. Evaluation designs that collect a graded score and then
threshold it are absorbing this instability silently; a binary or pairwise
formulation of the same underlying question is measurably more reproducible.

\subsection{Refusal is an evaluation failure mode, not an aside}

One judge declined a quarter of the relevance items while another, on identical
prompts, declined none. For anyone using an LLM judge over medical or otherwise
sensitive corpora this is a first-order problem: the missing verdicts are not
missing at random, and an evaluation harness that silently treats a decline as a
parse failure will report a format-adherence problem that does not exist.

It also breaks a convenience that most analyses assume. Because a rewording can
itself change whether a judge answers, refusal is an outcome of the manipulation
rather than a nuisance to be dropped, and conditioning on the answered subset is
selection on a post-treatment variable. We report worst-case bounds alongside the
conditioned estimate for this reason. Where refusals are common those bounds are
wide, and we would rather show a wide identified region than a precise number
that depends on an assumption nobody can check.

\subsection{On building a benchmark that can be wrong}

The defects we found in our own rebuild are, we think, the most transferable part
of this work. Each one passed every check that existed at the time.

A deterministic rotation assigning template pairs to items ran in lockstep with
the alternating emission of correct and incorrect answers, making the template
pair a perfect answer key --- and skewing two templates' label balance to
$75/25$, so a constant-answer judge would have appeared to have a large
template preference. Candidate responses reused across pairings carried
contradictory gold labels, so a judge reasoning correctly and consistently was
scored wrong on nearly half the preference items. A length-balancing step that
correctly removed a verbosity shortcut cut only one of the two buckets, and cut
it by vote margin, leaving the buckets matched on length and far apart on
annotator agreement.

None of these is detectable by inspecting labels, provenance, or aggregate
statistics. All three were found by asking what a heuristic with no understanding
would score, and by checking whether the artifact contained a rule that predicted
the answer. We would encourage that as routine practice: before reporting what
models score on a benchmark, report what nothing scores on it, and report it
with an interval and a support count rather than a bare number.

\subsection{Limitations}

The evidence base is three judges from one vendor family. The task ordering is
consistent across them, but that consistency is a within-family observation and
we do not claim it generalises; the harness supports fourteen judges across seven
providers and eleven are unrun. Two of the twelve judge--task cells report no
endpoint, their support having fallen below the pre-registered floor. One judge
rejects an explicit temperature parameter and therefore ran at a provider
default while the others ran at zero; this is recorded per call rather than
inferred, and that judge's cells should be read accordingly.

The design measures sensitivity within an item across a template pair, and cannot
compare templates to one another: no item appears under more than one pair, so
template identity is nested inside item identity and the two effects are not
separable. Our equivalence checks bound specific failure modes --- label space,
polarity, construct, non-triviality --- but nothing computable from the template
strings establishes that a competent reader maps two wordings onto the same
question. Semantic drift is reduced as an alternative explanation, not
eliminated.

Finally, the effect sizes are modest in absolute terms. A $\Delta\mathrm{JSS}$ of
$-0.04$ on factuality means a judge that reproduces itself on the identical
prompt loses four points of agreement under rewording. Whether that matters
depends on the decision being made downstream, and we have declared in advance
that we do not regard anything below $0.02$ as practically meaningful.

\section{Conclusion}

We set out to measure whether LLM judges return the same verdict when asked the
same question in different words, and to do so on a benchmark that could not be
passed without judging. Across three judges and four tasks, every measurable cell
shows a real loss of agreement under meaning-preserving rewording, beyond what
the judge's own decoding variance explains. The size of that loss depends more on
the evaluation task than on the model performing it.

The benchmark, the analysis plan fixed before the sweep, the raw judge outputs,
and the heuristic baselines that would beat the tasks are all released. So are
the defects we found in our own construction, and the measurements that exposed
them.


\subsubsection*{Reproducibility statement}
The dataset, loaders, analysis pipeline, pre-registered analysis plan, and every
raw judge output behind the reported numbers are released at
\url{https://github.com/rohithreddybc/judgeSense}. Each table is regenerated from
the committed raw outputs by a single script; no number in this paper is
transcribed by hand. Each split is content-hashed, and a 32-check audit gate runs
in continuous integration. Judge model identifiers, resolved decoding parameters
and the provider-reported model string are recorded on every call record, so
configuration drift is detectable after the fact rather than assumed absent.

\bibliographystyle{plainnat}          % TMLR: \bibliographystyle{tmlr}
\bibliography{references}

\end{document}
